% Sección del artículo — dueño: P1.
\section{Introducción}
Inspeccionar infraestructura con drones pequeños es barato y rápido, pero su navegación
depende casi por completo del GNSS (GPS). Esa señal se degrada justamente en lugares
frecuentes en la geografía peruana: en quebradas y cañones las paredes bloquean parte del
cielo y la señal rebota en la roca~\cite{groves2013}; debajo de puentes, en túneles y en
canales cubiertos no llega; y, por estar el Perú cerca del ecuador magnético, el centelleo
ionosférico ecuatorial puede hacer que los receptores pierdan la señal después del
atardecer~\cite{kintner2007}. Cuando esto ocurre, la función de «volver a casa» de los
drones comerciales queda sin referencia y el equipo deriva.

Los insectos resuelven un problema análogo con cerebros diminutos. Las abejas y hormigas
combinan la integración de la trayectoria con instantáneas visuales del entorno para
regresar a su nido~\cite{cartwright1983,wehner2003,srinivasan2011}, y las moscas esquivan
obstáculos con reflejos visuales rápidos y locales. Tomamos esa división del trabajo como
hipótesis de diseño: un módulo \emph{abeja} decide \emph{hacia dónde} ir y un módulo
\emph{mosca} decide \emph{cómo} moverse sin chocar.

Formulamos el problema como control secuencial bajo observación parcial. En cada paso el
dron recibe un parche egocéntrico de ocupación, un anillo de distancias, una panorámica de
baja resolución, la brújula, la odometría y, sólo mientras exista, el GPS; elige una de seis
acciones discretas. Una misión tiene éxito si visita todos los puntos de inspección y aterriza
a menos de 2\,m de la base sin chocar. El aprendizaje profundo es adecuado porque las
entradas son de alta dimensión (1\,024 celdas por cuadro), la decisión depende de la historia
reciente (obstáculos móviles, deriva) y el retorno visual exige comparar patrones
panorámicos de forma robusta al ruido, algo que los métodos analíticos como el vector
promedio de hitos (ALV)~\cite{lambrinos2000} hacen con errores de \nfbAlvErrorMediano{}.

Las contribuciones de este trabajo son:
\begin{itemize}
  \item Un simulador y protocolo de evaluación reproducibles para inspección con pérdida de
        GPS: seis condiciones, particiones por semilla sin fuga y un experto con información
        privilegiada que resuelve el \nfbExitoExperto{} de las misiones.
  \item Una arquitectura híbrida bioinspirada: \emph{mosca} (CNN+GRU) para el control local,
        \emph{abeja} (red siamesa 1D con relleno circular) para el retorno visual y un
        supervisor de seguridad independiente de las redes.
  \item Una comparación con cinco líneas base clásicas, ablaciones de arquitectura y de
        entrenamiento (imitación vs.\ DAgger), análisis de errores e interpretabilidad
        (Grad-CAM e Integrated Gradients).
  \item Una validación del sensor panorámico 2.5D frente a propiedades conocidas de los
        panoramas reales, una ablación de la red de la abeja en la partición de prueba y
        pruebas de significancia pareadas para cada comparación.
\end{itemize}
